% Part: sets-functions-relations
% Chapter: relations
% Section: orders

\documentclass[../../../include/open-logic-section]{subfiles}

\begin{document}

\olfileid[tr]{sfr}{rel}{ord}
\olsection{Sıralamalar}

\begin{explain}
Karşılaştırmalarımızın birçoğunda, belirli bir bakımdan bazı nesneleri başka
nesnelerden ``küçük'', onlara ``eşit'' ya da onlardan ``büyük'' diye niteleriz.
Bu karşılaştırmalar \emph{sıralama} bağıntılarını içerir. Ancak sıralama
bağıntılarının farklı türleri vardır. Örneğin bazıları herhangi iki nesnenin
karşılaştırılabilir olmasını gerektirirken bazıları gerektirmez. Bazıları
özdeşliği içerir ($\le$ gibi), bazılarıysa dışlar ($<$ gibi). Burada bir
sınıflandırma yapmak yararlı olacaktır.
\end{explain}

\begin{defn}[Ön sıralama]
Hem yansımalı hem de geçişli olan bir bağıntıya \emph{ön sıralama} denir.
\end{defn}

\begin{defn}[Kısmi sıralama]
Aynı zamanda ters simetrik olan bir ön sıralamaya \emph{kısmi sıralama} denir.
\end{defn}

\begin{defn}[Doğrusal sıralama]\ollabel{def:linearorder}
Aynı zamanda bağlantılı olan bir kısmi sıralamaya \emph{tam sıralama} veya
\emph{doğrusal sıralama} denir.
\end{defn}

\begin{ex}
Her doğrusal sıralama aynı zamanda kısmi sıralama, her kısmi sıralama da aynı
zamanda ön sıralamadır; fakat tersleri doğru değildir. $A$ üzerindeki evrensel
bağıntı, yansımalı ve geçişli olduğu için bir ön sıralamadır. Ancak $A$'nın
birden fazla elemanı varsa evrensel bağıntı ters simetrik değildir ve dolayısıyla
bir kısmi sıralama değildir.
\end{ex}

\begin{ex}
$\Bin^*$ üzerinde $\preccurlyeq$ ile göstereceğimiz \emph{daha uzun değildir}
bağıntısını ele alalım: $x \preccurlyeq y$ ancak ve ancak $\len{x} \le \len{y}$. Bu, yansımalı
ve geçişli bir ön sıralamadır; hatta bağlantılıdır, fakat ters simetrik olmadığı
için bir kısmi sıralama değildir. Örneğin $01 \preccurlyeq 10$ ve
$10 \preccurlyeq 01$ olduğu hâlde $01 \neq 10$'dur.
\end{ex}

\begin{ex}
Önemli bir kısmi sıralama, bir kümeler kümesi üzerindeki $\subseteq$
bağıntısıdır. Bu genel olarak doğrusal bir sıralama değildir. Çünkü $a \neq b$
ve $\Pow{\{a, b\}}=\{\emptyset, \{a\}, \{b\}, \{a,b\}\}$ kümesini ele
aldığımızda $\{a\} \nsubseteq \{b\}$, $\{a\} \neq \{b\}$ ve
$\{b\} \nsubseteq \{a\}$ olduğunu görürüz.
\end{ex}

\begin{ex}
\emph{Kalansız bölünebilme} bağıntısı bize doğrusal olmayan bir kısmi sıralama
verir. $n$ ve~$m$ tam sayıları için $n \mid m$ yazımı, $n$'nin $m$'yi kalansız
böldüğü; yani $m=kn$ olacak biçimde bir $k$ tam sayısının bulunduğu anlamına
gelsin. Bu bağıntı $\Nat$ üzerinde bir kısmi sıralamadır fakat doğrusal sıralama
değildir; örneğin $2 \nmid 3$ ve $3 \nmid 2$. Bölünebilme, $\Int$ üzerinde bir
bağıntı olarak ele alındığında ters simetrik olmadığı için yalnızca bir ön
sıralamadır: $1 \mid -1$ ve $-1 \mid 1$ olduğu hâlde $1 \neq -1$'dir.
\end{ex}

\begin{ex}
$A^*$ dizileri kümesi üzerindeki \emph{uzantı} bağıntısı şöyledir:
$s \sqsubseteq s'$ ancak ve ancak $s=\emptyseq$ (boş dizi), $s=s'$ veya
$s=\tuple{s_1, \dots, s_n}$ ve
$s'=\tuple{s_1, \dots, s_n, s_{n+1}, \dots, s_m}$. $s \sqsubseteq s'$ ise
$s$'ye ayrıca $s'$ dizisinin bir \emph{başlangıç kesiti} deriz. $A^*$
üzerindeki uzantı bağıntısı kısmi bir sıralamadır, fakat doğrusal bir sıralama
değildir. Örneğin $a \neq b$ ise $ab \not\sqsubseteq ba$ ve
$ba \not\sqsubseteq ab$ olur.
\end{ex}

\begin{defn}[Sıkı sıralama]
Yansımasız, asimetrik ve geçişli olan bir bağıntıya \emph{sıkı sıralama} denir.
\end{defn}

\begin{defn}[Sıkı doğrusal sıralama]\ollabel{def:strictlinearorder}
Aynı zamanda bağlantılı olan bir sıkı sıralamaya \emph{sıkı tam sıralama} veya
\emph{sıkı doğrusal sıralama} denir.
\end{defn}

\begin{ex}
$\le$, sıkı doğrusal sıralama~$<$'ye karşılık gelen doğrusal sıralamadır.
$\subseteq$ ise sıkı sıralama~$\subsetneq$'ye karşılık gelen kısmi sıralamadır.
\end{ex}

$A$ üzerindeki herhangi bir sıkı sıralama $R$'ye köşegen $\Id{A}$, yani bütün
$\tuple{x,x}$ ikilileri eklenerek bir kısmi sıralama elde edilebilir. (Buna
$R$'nin \emph{yansımalı kapanışı} denir.) Tersine, bir kısmi sıralamadan
başlayıp $\Id{A}$ çıkarılarak sıkı bir sıralama elde edilebilir. Sonraki iki
sonuç bunu kesin olarak ifade eder.

\begin{prop}\ollabel{prop:stricttopartial}
$R$, $A$ üzerinde sıkı bir sıralamaysa $R^+=R\cup\Id{A}$ bir kısmi sıralamadır.
Ayrıca $R$ sıkı bir doğrusal sıralamaysa $R^+$ doğrusal bir sıralamadır.
\end{prop}

\begin{proof}
$R$'nin sıkı bir sıralama olduğunu; yani $R \subseteq A^2$ olup $R$'nin
yansımasız, asimetrik ve geçişli olduğunu varsayalım. $R^+=R\cup\Id{A}$ olsun.
$R^+$'nın yansımalı, ters simetrik ve geçişli olduğunu göstermeliyiz.

Her $x \in A$ için $\tuple{x,x}\in\Id{A}\subseteq R^+$ olduğundan $R^+$
açıkça yansımalıdır.

$R^+$'nın ters simetrik olduğunu göstermek için, olmayana ergi yoluyla $R^+xy$
ve $R^+yx$ geçerli olduğu hâlde $x \neq y$ olduğunu varsayalım.
$\tuple{x,y}\in R\cup\Id{A}$, fakat $\tuple{x,y}\notin\Id{A}$ olduğundan
$\tuple{x,y}\in R$; yani $Rxy$ olmalıdır. Benzer biçimde $Ryx$ olur. Ancak bu,
$R$'nin asimetrik olduğu varsayımıyla çelişir.

Geçişliliği göstermek için $R^+xy$ ve $R^+yz$ olduğunu varsayalım. Hem
$\tuple{x,y}\in R$ hem de $\tuple{y,z}\in R$ ise, $R$ geçişli olduğundan
$\tuple{x,z}\in R$ olur. Aksi hâlde ya $\tuple{x,y}\in\Id{A}$, yani $x=y$;
ya da $\tuple{y,z}\in\Id{A}$, yani $y=z$ olur. İlk durumda varsayım gereği
$R^+yz$ geçerlidir ve $x=y$ olduğundan $R^+xz$ elde edilir. İkinci durumda da
benzer biçimde aynı sonuç çıkar. Her iki durumda $R^+xz$ olduğuna göre $R^+$
geçişlidir.

``Ayrıca'' bölümüne gelince, $R$'nin bağlantılı da olduğunu varsayalım.
Dolayısıyla $A$'nın bütün farklı $x,y$ elemanları için ya $Rxy$ ya da $Ryx$; yani ya
$\tuple{x,y}\in R$ ya da $\tuple{y,x}\in R$ olur. $R\subseteq R^+$ olduğundan
bu $R^+$ için de geçerliliğini korur; dolayısıyla $R^+$ da bağlantılıdır.
\end{proof}

\begin{prop}\ollabel{prop:partialtostrict}
$R$, $A$ üzerinde bir kısmi sıralamaysa $R^-=R\setminus\Id{A}$ sıkı bir
sıralamadır. Ayrıca $R$ doğrusal bir sıralamaysa $R^-$ sıkı bir doğrusal
sıralamadır.
\end{prop}

\begin{proof}
Bu, alıştırma olarak bırakılmıştır.
\end{proof}

\begin{prob}
\olref[sfr][rel][ord]{prop:partialtostrict} için bir ispat verin.
\end{prob}

Aşağıdaki basit sonuç, sıkı doğrusal sıralamaların kaplamsallık benzeri bir
özelliği sağladığını gösterir:

\begin{prop}\ollabel{prop:extensionality-strictlinearorders}
$<$, $A$ üzerinde sıkı bir doğrusal sıralamaysa
\[
  (\forall a, b \in A)((\forall x \in A)(x < a \liff x < b) \lif a = b).
\]
\end{prop}

\begin{proof}
$(\forall x \in A)(x<a\liff x<b)$ olduğunu varsayalım. $a<b$ olsaydı
$a<a$ olurdu; bu da $<$'nin yansımasız olmasıyla çelişirdi. Öyleyse
$a\nless b$. Tam olarak aynı biçimde $b\nless a$. $<$ bağlantılı olduğundan
$a=b$ sonucuna varırız.
\end{proof}

\end{document}
